Drehmoment $M = F\cdot d$, mit dem Hebelarm $d$ als Abstand von D zur Wirkungslinie. Der Gegenuhrzeigersinn zählt positiv.

\begin{abcliste}
\abc
$F_1$: Bei der Steigung $3:4$ ist die schräge Strecke 5 Kästchen lang (3-4-5-Dreieck). Die Komponenten sind
\begin{align}
 F_{1x} &= \FaxF = \Fax, & F_{1y} &= \FayF = \Fay
\end{align}
Beide greifen im Angriffspunkt an, $x_1 = \xaO$ rechts von und $y_1 = \yaO$ über D. Die senkrechte Komponente hat den Hebelarm $x_1$ und dreht im Gegenuhrzeigersinn, die waagrechte hat den Hebelarm $y_1$ und dreht (nach rechts, oberhalb von D) im Uhrzeigersinn:
\begin{align}
 M_1 &= \MaF = \Fay\cdot\xa - \Fax\cdot\ya = \May - \Max \\
   &= \Ma \approx \result{\MaP{0}{2}} \quad\text{(Gegenuhrzeigersinn)}
\end{align}

$F_2$: Ihre Wirkungslinie ist die waagrechte Gerade durch D. Der Hebelarm ist null:
\[
 M_2 = \result{0}
\]

$F_3$: senkrecht nach oben, $x_3 = \xcO$ links von D, also im Uhrzeigersinn:
\begin{align}
 M_3 &= \McF = \Fc\cdot\xc = \Mc \approx \result{\McP{0}{2}} \quad\text{(Uhrzeigersinn)}
\end{align}

$F_4$: waagrecht nach links, Wirkungslinie $y_4 = \ydO$ unter D, also im Uhrzeigersinn:
\begin{align}
 M_4 &= \MdF = \Fd\cdot\yd = \Md \approx \result{\MdP{0}{2}} \quad\text{(Uhrzeigersinn)}
\end{align}

\abc
\[
 \result{M_2 < M_1 < M_3 < M_4}
\]
Die Kräfte selbst sind genau umgekehrt geordnet, $F_4 < F_3 < F_2 < F_1$: Die grösste Kraft $F_1$ hat das zweitkleinste Drehmoment, weil sich ihre beiden Komponenten fast aufheben, und $F_2$ hat gar keines. Die kleinste Kraft $F_4$ hat den längsten Hebelarm und das grösste Drehmoment.

\abc
\begin{align}
 d_1 &= \daF = \frac{\Ma}{\Fa} = \da \approx \result{\daP{-2}{2}}
\end{align}
Der Abstand von D zum Angriffspunkt ist $\sqrt{20^2+10^2}\,\si{\cm}\approx\SI{22}{\cm}$, mehr als fünfmal so viel. Der Hebelarm ist das Lot von D auf die verlängerte Wirkungslinie und ist immer höchstens so lang wie der Abstand zum Angriffspunkt.

\abc
Im Uhrzeigersinn drehen $M_3$ und $M_4$, im Gegenuhrzeigersinn nur $M_1$:
\begin{align}
 M_\mathrm{res} &= \MrF = \Mc + \Md - \Ma = \Mr \approx \result{\MrP{0}{2}}
\end{align}
im Uhrzeigersinn. Die Platte beginnt sich \result{\text{im Uhrzeigersinn}} zu drehen.
\end{abcliste}
